\section{Introduction}
In 1754 a coach advertised that it would carry passengers from London to Edinburgh in ten days in summer and twelve in winter \citep{harper1922}. In 1837 the Edinburgh mail coach was timed at 42 hours and 23 minutes \citep{harper1903}. By the early 1900s the Flying Scotsman ran from King's Cross to Edinburgh in eight hours and a quarter \citep{wiki_fs}, and today the fastest train is timed at four hours \citep{visitlondon}. These are the familiar milestones of the transport revolution, and they have been told many times as falling journey times between pairs of famous places. What they do not show is the whole map: how much of the country lay within a day of the capital, and where the edge of that day fell, in each era.

This paper draws that map. It asks a simple question: leaving Charing Cross at dawn, how much of Great Britain could a traveller reach within twelve hours of travel in 1750, in the 1830s, in 1850, in 1910 and today? To answer it I build a transparent multimodal network model on a 2 km grid. Its road layer for 1750 and the 1830s is the turnpike and main road map of England and Wales digitised by \citet{rosevear2023tp}, with each section coded by whether it was a turnpike in that decade; its coach speeds are calibrated, one parameter per era, on documented journey times, including {{coach.n_termini_1836|int}} London destinations with scheduled times in 1836 \citep{bates1969,rosevear2023coach}. Its rail layer for 1850 is the 1851 network digitised by the Cambridge Group \citep{satchell2023}, for 1910 a reconstruction from that network, today's lines and the stations that closed after 1910, and today's network from OpenStreetMap. Train speeds come from Leunig's (\citeyear{leunig2006}) estimates of average speeds on English and Welsh railways, and everything is checked against documented journeys that were not used to set the speeds.

The answer is a sequence of five maps (Figure~\ref{fig:maps}). In 1750 a day's coaching from London covered {{A.1750.km2|int}} km$^2$, {{A.1750.share|f0}}\% of Great Britain's land, a patch reaching roughly {{A.1750.reach_km|f0}} km from Charing Cross. On the faster coaches and better roads of the 1830s it covered {{A.1830.share|f0}}\%. The railway of 1850 pushed the edge of the day to {{A.1850.share|f0}}\%, the network of 1910 to {{A.1910.share|f0}}\%, and today {{A.today.share|f0}}\% of the land can be reached within twelve hours by car or train. These numbers are not equally firm. The 1750 model misses its documented journeys by more than the limit set in advance, so the 1750 map is a rough sketch; the railway models use average trains and are much slower than the documented expresses, so the 1850 and 1910 areas are lower bounds (Section~\ref{sec:results}). Between 1750 and today the area within a day of London grew {{ratio.today_1750|f0}}-fold. The most rapid shrinking of the country happened twice: between 1750 and the 1830s, on the roads, and between the 1830s and 1850, when the first railways more than {{RATIO_1850_TXT}} the day's reach in about fifteen years.

The paper makes three contributions. First, it turns the well known point-to-point story of time-space convergence \citep{janelle1968,janelle1969} into a continuous, national surface for five dates, built from published historical GIS of roads and railways rather than from straight lines or a handful of routes. Second, it calibrates and tests that surface openly: every speed is either fitted to cited timetables or taken from a published estimate, and every era's model is compared with journeys it was not fitted to, so readers can see how far to trust each map. The 1830s model, fitted on the 1836 coach timetables, reproduces the 1837 mail-coach timetable to within {{valid.1830_mail1837.mdape|f0}}\% (median absolute error); the 1750 model, built on a dozen cited journeys, fails the error threshold set in advance and is a rough sketch; and the railway models, checked against only four and six documented express journeys, are much slower than those expresses, so their areas are lower bounds. Third, the derived grids, isochrones and the compiled journey times are released for reuse.

Section~\ref{sec:lit} reviews the literature on time-space convergence, pre-railway road transport and railway speeds. Section~\ref{sec:data} describes the data, Section~\ref{sec:model} the model, its calibration and the analysis plan, Section~\ref{sec:results} the results, Section~\ref{sec:robust} the checks, and Section~\ref{sec:disc} the limitations.

\section{Background and literature}\label{sec:lit}
\subsection{Time-space convergence and accessibility}
Janelle introduced the idea of time-space convergence: as transport improves, the travel time between places falls, so that in time, if not in space, places move closer together \citep{janelle1968}. He measured it as the average rate at which travel time between a pair of places falls, in minutes per year, and linked it to the reorganisation of activities around the places that gain most \citep{janelle1969}. The concept has since been carried into time-space maps, which redraw a territory so that distances on the page are proportional to travel times \citep{spiekermann1994}, and into long-run measures of accessibility, such as the Swiss reconstruction since 1850 by \citet{axhausen2011}, which found road accessibility overtaking rail around 1930. Cultural histories describe the same change as a ``shrinking'' of space and an industrialisation of time \citep{schivelbusch1986}. Historical work on London itself shows how rail access and population density developed together \citep{levinson2008}. These studies share a need for travel times that are comparable across eras and across space. The present paper supplies them for Britain from a single model applied to five network states.

\subsection{Roads and coaches before the railway}
The literature on the turnpike era is rich in quantitative evidence. Turnpike trusts, local bodies empowered by Acts of Parliament to levy tolls and improve a stretch of road, are the subject of classic studies \citep{albert1972,pawson1977}; by 1838, on the map used here, nearly 23,000 miles of road were under active turnpike trusts \citep{rosevear2023tp}. Historical atlases map coach journey times from London around 1750 and in the 1830s \citep{pope1989}. \citet{bogart2005} assembled journey times from the coaching literature and showed that average passenger journey speeds rose from {{BOG1750}} miles per hour in the 1750s to {{BOG1820}} in the 1820s, measuring distance over total elapsed journey time, and linked the gains to turnpike investment: the London to Birmingham journey fell from 65 hours in 1705 to 40 hours in 1740 and 25 hours in 1770 as trusts took over the route. \citet{gerhold2014}, using newspaper advertisements from 1653 to 1840, dated the acceleration of London coaches to two periods, the 1760s to 1780s and the 1810s to 1820s, and attributed it to turnpikes, better vehicles and better horses. \citet{gerhold2005} documents the carriers and coaches of the pre-turnpike era.

The Cambridge Group's transport project has turned this evidence into historical GIS. The turnpike map used here records every turnpike and main road of England and Wales from about 1680 to 1850, built from Cary's maps and the Acts, with the decade in which each section was a turnpike \citep{rosevear2023tp}. A companion file maps the routes and scheduled times of London coach services in 1819 (from Cary's \textit{New Itinerary}) and 1836 \citep{bates1969,rosevear2023coach}. These data have been used to study the geography of coaching \citep{rosevear2018} and, with assumed average speeds per mode, to estimate transport costs and their fall between 1680 and 1830 \citep{alvarez2017}, who assume 2 miles per hour for passengers by road in 1680 and 8.7 mph in 1830 (the average for journeys of 80 miles and more in Gerhold's sample). The present paper uses the same road map but fits the speed to scheduled times rather than assuming it, and extends the comparison to the railway era and today.

Older narrative sources record many individual journeys. Smiles's history of roads and travelling \citep{smiles1867} and Harper's road books \citep{harper1899,harper1902,harper1903,harper1922} quote advertisements, diaries and time-bills. Jackman's survey remains the fullest general account \citep{jackman1916}. These sources are uneven and sometimes repeat each other, and a stated duration in days leaves the hours on the road unknown. They are used here for what they are: a set of documented journeys, each cited, against which a model can be checked.

\subsection{Railway speeds}
The railway network of England, Wales and Scotland has been reconstructed as a time-dynamic GIS from Cobb's atlas by the Cambridge Group and the University of Lleida, with snapshots for 1851, 1861 and 1881 released publicly \citep{marti2017,satchell2023}. It has been used to measure how quickly railways reached the population and whether stations brought growth \citep{gregory2010,bogart2022}, and set in a European context \citep{marti2013}. For speeds, \citet{leunig2006} transcribed timetables for important and minor routes in England and Wales in 1850, 1870, 1887 and 1910 and estimated average in-train speeds, measured on a straight-line basis, of {{speeds.1850.leunig_important_mph|f1}} mph on important journeys and {{speeds.1850.leunig_minor_mph|f1}} mph on minor ones in 1850, rising to {{speeds.1910.leunig_important_mph|f1}} and {{speeds.1910.leunig_minor_mph|f1}} mph in 1910 \citep[Table~2]{leunig2006}. Those averages are the train speeds of this model.

\subsection{Contribution}
What is missing is a single, testable surface of travel time across the whole of Great Britain that can be compared across the coach and rail eras. Point-to-point histories cannot say how much of the country lay within a day of London; maps such as Pawson's give isochrones for coach journeys but not a common method that carries through the railway age; and economic studies that build network models usually assume speeds rather than test them against timetables. This paper builds one model, calibrates one speed per era where data allow, takes published speeds where they do not, and reports how well each era's map reproduces journeys it has not seen.

\section{Data}\label{sec:data}
\subsection{Roads, 1750 and the 1830s}
The road layer for England and Wales is the ``turnpikes and main'' shapefile of \citet{rosevear2023tp}: {{TP_N}} line parts (23,556 features before multi-part features are split) and about 41,000 miles of road, including every turnpike, the ``other main roads'' of Cary's maps, link roads, bridges and ferries. For each decade from 1750 to 1850 a field records whether the section was a turnpike at that date; roads that were not yet built (new alignments) are coded separately and dropped from earlier networks. In the 1750 network a section is either a turnpike in 1750 or an existing main road that was not; the same for 1830 (and 1820 for the 1819 check, 1838 for the feeder roads of 1850 and 1910). The map does not cover Scotland. For Scotland I use today's trunk and primary roads from OpenStreetMap as a proxy for its main roads, at the non-turnpike speed in 1750 and the turnpike speed in the 1830s. This matters only for the far north, which lies well beyond a day from London in both coach eras.

\subsection{Coach timetables, 1819 and 1836}
The London coaches file \citep{rosevear2023coach} maps {{CO_ROUTES}} coach routes that start in London, each drawn on the turnpike map, with scheduled journey times where the source gives one. Leaving out short stages whose terminus is less than 15 km from Charing Cross, {{coach.n_services_1836|int}} timed services in 1836 reach {{coach.n_termini_1836|int}} distinct termini, and {{coach.n_services_1819|int}} services in 1819 reach {{coach.n_termini_1819|int}}. The median scheduled speed of all timed 1836 services, measured along the mapped route, is {{coach.median_mph_1836|f1}} mph, and {{coach.median_mph_80plus_1836|f1}} mph for routes of 80 miles or more, close to the 8.7 mph that \citet{alvarez2017} take from Gerhold. For calibration I keep the fastest service to each terminus, because the question is how far a traveller could get, and take the far end of the route as the terminus point.

\subsection{Railways, 1850 and 1910}
The 1850 rail layer is the 1851 England, Wales and Scotland rail lines snapshot \citep{satchell2023}: {{R51_N}} polylines, 6,336 miles of line open for passengers or freight during 1851 (the file does not distinguish the two). There is no public snapshot for 1910: the full time-dynamic file is restricted to the project teams \citep{marti2017}. I therefore reconstruct the 1910 network from three parts: the 1851 lines (nearly all still open in 1910); today's railway lines from OpenStreetMap; and lines that have since closed, inferred from the {{NFAR}} stations in a list of closed British stations that closed after 1910 and lie more than 2 km from either of the first two sets. These stations are joined by a minimum spanning tree with links of at most 15 km, and each group to the nearest line, adding {{INF_KM}} km of inferred line. The inferred lines are straight segments between stations, so their lengths are too short where real lines curved, and a few will join stations that were never connected. Lines opened after 1910 and still open (the Channel Tunnel Rail Link, for example) are not removed. Rail sections are classed as main lines when at least half their length lies within 250 m of a line that OpenStreetMap shows at zoom level 8 (a proxy for today's main lines) and as branch lines otherwise.

Stations determine where a traveller can board. In 1850 they are the open stations whose Wikidata opening year is 1850 or earlier, plus closed stations within 500 m of an 1851 line; in 1910, open stations opened by 1910 plus every station closed after 1910; today, the stations open today. A rail vertex within 1.5 km of a station can be boarded.

\subsection{Today}
Today's road layer is the OpenStreetMap network of motorways, trunk, primary and secondary roads, read from OpenFreeMap vector tiles at zoom level 10 (OSM positional accuracy is high on main roads; \citealp{haklay2010}). Today's rail layer is the OpenStreetMap railway network (standard gauge lines, not sidings, trams or the Underground), with today's open stations.

\subsection{Documented journeys}
For each era I compiled documented journey times from London with a source for each (Appendix Table~\ref{tab:docs}). For 1750, the twelve journeys between 1739 and 1764 that I could find and cite, to nine towns: from \citet{smiles1867}, \citet{bogart2005}, Harper's road books and a summary of Gerhold's York services of 1764 \citep{gnrsite}. Journeys given in days are converted at twelve travelling hours a day, which is also the length of the day modelled; the sensitivity to ten and fourteen hours is reported. The thirty to sixty journeys per era that a firm calibration would want could not be found in sources reachable for this study, and this is the weakest part of the evidence. For the 1830s, the 1836 timetables above for calibration and, as an out-of-sample test, the list of {{valid.1830_mail1837.n|int}} mail coaches leaving London nightly in 1837 printed by \citet{harper1903}. For the railway eras, the fastest documented times I could cite: to Liverpool in 1851 \citep{sidney1851}, Aberdeen in 1855 \citep{wiki_race}, Edinburgh in 1862 \citep{wiki_fs} and Birmingham on the first timetable of 1838 \citep{tring} for 1850; Edinburgh, Glasgow, Plymouth, Penzance, Bristol and Aberdeen between 1896 and 1904 for 1910 \citep{wiki_fs,harper1903,wiki_cre,wiki_cornishman,wiki_race}; and the fastest times from London to twelve cities stated by \citet{visitlondon} for today.

\section{Model and design}\label{sec:model}
\subsection{Graph}
Great Britain is divided into 2 km cells on the British National Grid; a cell is land if its centre lies inside a local authority of England, Wales or Scotland, giving {{LAND_CELLS}} cells and {{land_km2|int}} km$^2$. Each era's graph has three kinds of edge. Walking edges join each land cell to its eight neighbours at 4.8 km/h (3 mph); today these edges stand for minor roads at 30 km/h. Network edges follow each road or rail line, resampled every 300 m, at the speed of that line; vertices of different lines within 40 m of each other are joined, which connects junctions and the seams of the vector tiles. Access edges join each network vertex to the cell that holds it, at walking speed for the distance between them, with a wait of 30 minutes on boarding a train (and, in 1850 and 1910, a feeder coach); alighting is free. Ferries in the turnpike map carry a 30 minute wait. Travel times from the cell holding Charing Cross are found with Dijkstra's algorithm \citep{dijkstra1959}; the time to a town is the smallest time to any node within 2.5 km plus the walk from that node. The model has no timetables: it assumes a service is there when the traveller is, beyond the fixed waits.

\subsection{Speeds and calibration}
In the coach eras the free parameter is the speed of a coach on a turnpike, $v_T$, with existing main roads that were not turnpikes at $r = 0.6$ of that speed (a stated assumption, varied from 0.4 to 1.0). For the 1830s, $v_T$ is chosen on a grid of 0.5 km/h steps to minimise the sum of squared log errors between modelled and scheduled times for the fastest 1836 service to each of the {{coach.n_termini_1836|int}} termini; for 1750, on a grid of 0.25 km/h steps against the twelve cited journeys. The same fit on the 1819 times, using the turnpikes of 1820, checks that the fitted speeds follow the trajectory in the literature. In 1850 and 1910 feeder coaches run on the 1838 turnpike network at the 1830s speeds.

Train speeds in 1850 and 1910 come from \citet[Table~2]{leunig2006}. His averages are in-train speeds measured along the straight line between the ends of a journey; I convert them to speeds along the track by multiplying by the median circuity of each era's rail network (network distance over straight-line distance) from Charing Cross to the 30 largest towns of England and Wales on the town list that lie within 5 km of each network ({{speeds.1850.n_towns}} in 1850, {{speeds.1910.n_towns}} in 1910): {{speeds.1850.circuity|f2}} in 1850 and {{speeds.1910.circuity|f2}} in 1910. Main lines run at Leunig's important-journey speed, {{RAIL1850_KMH}} km/h in 1850 and {{RAIL1910_KMH}} km/h in 1910, and branch lines at the ratio of his minor to important speeds ({{speeds.1850.branch_ratio|f2}} and {{speeds.1910.branch_ratio|f2}}). These speeds describe an average train, not the fastest express. Today's main-line speed is fitted, as in the coach eras, to the twelve fastest times from London, with branch lines at half that speed; today's road speeds are fixed at 100 km/h on motorways, 80 on trunk roads, 60 on primary and 45 on secondary roads, below the 2024 average free-flow car speeds of 68 mph on motorways and 51 mph on rural single carriageways \citep{dft2025} to allow for junctions, towns and traffic. A traveller today takes whichever of car or train is faster.

\subsection{Analysis plan}
The analysis was pre-specified in a written plan, kept in the study notes, before any travel time was computed: the twelve-hour day, the 2 km grid, the graph and its fixed parameters, one fitted speed per era, leave-one-destination-out validation, the out-of-sample tests, the headline outputs and the sensitivity checks below. One change was made before any calibration or validation was run, and is recorded in the notes: too few documented railway times could be found to fit train speeds for 1850 and 1910, so those come from Leunig's published averages, with a fit to the documented fastest trains kept as a sensitivity check. The plan named in advance what would count against the model: a leave-one-out median absolute error above 25\% in an era, or a mean log error beyond $\pm0.1$.

\subsection{Validation and uncertainty}
For the fitted eras, validation is leave one destination out: the speed is refitted without the destination and its time predicted. I report the number of journeys, the median absolute percentage error (MdAPE), the share within 20\% and the mean log error (bias). For 1850 and 1910 every documented journey is out of sample. On the maps the uncertainty of each era is shown by redrawing the twelve-hour line with all speeds raised and lowered by that era's median absolute error.

\section{Results}\label{sec:results}
\subsection{Calibrated speeds}
{{RES_SPEEDS}}

\subsection{How well the model reproduces documented journeys}
{{RES_VALID}}

\subsection{A day from London}
{{RES_MAPS}}

\subsection{Towns and the road north}
{{RES_TOWNS}}

\section{Robustness}\label{sec:robust}
{{RES_ROBUST}}

\section{Discussion}\label{sec:disc}
{{DISCUSSION}}

\section{Conclusion}\label{sec:conc}
{{CONCLUSION}}
